\section{Reasoning：语言模型智能的基石}

\subsection{什么是 Reasoning}

\begin{importantbox}{Reasoning 的核心定义}
Reasoning（推理）是一种"无中生有"的能力——从已知且有限的前提（Premise）出发，沿着推理链推导出更多、更丰富、更有深度的结论。
\end{importantbox}

经典的演绎推理三段论就是一个典型的例子：
\begin{itemize}
    \item 大前提：所有人都会死
    \item 小前提：苏格拉底是人
    \item 结论：苏格拉底也会死
\end{itemize}

模型真正的智能就是通过 Reasoning 来体现的。当我们对模型的回答感到 impressed 或 disappointed 时，本质上都是对其 Reasoning 能力的一种感知——超出预期令人惊叹，低于预期则令人失望。

\subsection{Reasoning 与 Agent 能力的关系}

\begin{knowledgebox}{Reasoning 是一切高级能力的基础}
如果模型没有 Reasoning 的能力：
\begin{itemize}
    \item 就没有 \textbf{Planning}：因为 Planning 需要分解任务、Think Ahead、想多步
    \item 就没有 \textbf{Reflection / Critique}：因为反思需要从已有前提产生大量新结论
    \item 就没有有效的 \textbf{Agentic} 行为：因为 Agent 需要基于推理来决策下一步动作
\end{itemize}
\end{knowledgebox}

\subsection{两类经典推理方式}

\begin{enumerate}
    \item \textbf{演绎推理（Deduction）}：前向推理，是确定性推理。从大前提和小前提出发推导出必然结论。
    \item \textbf{归纳推理（Induction）}：是或然性推理。从有限的观察中归纳出一般性结论，前提越多结论越自信，但不能保证绝对正确（如"天鹅都是白色的"）。
\end{enumerate}

\subsection{Reasoning Model：隐式的 Long Reasoning}

当前的 Reasoning Model（如 DeepSeek R1、OpenAI O 系列）追求的是一种\textbf{隐式的 Long Reasoning}：
\begin{itemize}
    \item 复杂问题 $\rightarrow$ 模型自发产生长的思考过程
    \item 简单问题 $\rightarrow$ 产生短的思考过程
    \item 这是在训练阶段让模型自发学会的能力
\end{itemize}

\subsection{本章小结}

Reasoning 是语言模型一切高级能力的基石。无论是 Planning、Reflection 还是 Agent 行为，都依赖于模型从有限前提中推导出新结论的能力。现代 Reasoning Model 正在通过 RL 训练来增强这种隐式的长链推理能力。


